\begin{figure}[H]
\centering
\begin{tikzpicture}[>=Stealth,scale=1.0]
  \draw[->] (-4,0) -- (7.25,0) node[right] {eje óptico};
  \filldraw[fill=cyan!15,draw=blue!60,thick]
    (0,-2.2) .. controls (-0.85,-1.1) and (-0.85,1.1) .. (0,2.2)
    .. controls (0.85,1.1) and (0.85,-1.1) .. cycle;
  \node[above] at (0,2.25) {objetivo};
  \foreach \y in {-1.25,1.25}{
    \draw[black!55,thick] (-3.7,\y) -- (0,\y);
    \draw[blue,very thick] (0,\y) -- (3.0,0);
    \draw[red,very thick] (0,\y) -- (4.65,0);
  }
  \fill[blue] (3,0) circle (2pt);
  \fill[red] (4.65,0) circle (2pt);
  \fill[yellow!65,opacity=.3] (3.82,0) ellipse (0.85 and 0.22);
  \draw[dashed,blue] (3,-0.4) -- (3,0.4);
  \draw[dashed,red] (4.65,-0.4) -- (4.65,0.4);
  \draw[<->,thick] (3,1.55) -- (4.65,1.55)
    node[midway,above,align=center] {intervalo focal\\cromático};
  \draw[<->,blue] (0,-1.75) -- (3,-1.75) node[midway,below] {$f_{\text{azul}}$};
  \draw[<->,red] (0,-2.2) -- (4.65,-2.2) node[midway,below] {$f_{\text{rojo}}$};
  \draw[blue,opacity=.4,line width=3pt] (5.8,-0.55) -- (5.8,0.55);
  \draw[red,opacity=.4,line width=3pt] (6.08,-0.55) -- (6.08,0.55);
  \node[align=left,anchor=west] at (6.25,0.75) {imagen con\\bordes coloreados};
\end{tikzpicture}
\caption{Aberración cromática longitudinal en el objetivo de un refractor.}
\end{figure}